% !TeX root = ../rvmonitor.tex

\section{Terms and Terminology}

%\subsection{Propositional and Higher-Order Logic}

\subsection{Temporal Logic}
Temporal logic is an extension of propositional logic integrating temporal relationships.
There exist different types of temporal logics which provide specific extension or are tailored for certain applications.
There exist temporal logics, i.e., the semantics, which can cover specifications forward, backward, forward and backward, over real-valued signals/events or only boolean, and also temporal logics which extend to regular expressions or consider first-order logic operators.´
The basic temporal logic is called linear temporal logic (LTL).
It extends propositional logic with operators specifying temporal operations such that something has to hold at the next time step or that something must hold until something else is valid.
A further 

\subsection{Runtime Verification}
We first want to introduce some important terms.
An event is a tuple $(\tau_i, \Gamma_i)$, where $i \in \mathcal{N}_0$ is a non-negative index ,$\tau_i \in \mathcal{T}$ is a time-stamp, and $ \Gamma_i \subseteq \mathcal{AP}$ is a subset of atomic predicates.
A stream is an infinite sequence of events, where a finite prefix of stream is denoted as trace.

Runtime Verification (RV) is an approach for monitoring a system.
In RV, a system is verified by analyzing system events against a specification.
Usually the specification is given in a temporal logic, e.g., LTL, MTL, MFOTL, or STL.
RV is highly-related to model checking but model checking verifies whether a system description complies with the specification whereas RV only uses a stream or trace.

RV is distinguished between online and offline monitoring. 
Online monitoring processes events at runtime while offline monitoring processes them after the system has stopped.
Speaking more general, offline monitoring assumes that all events are given or could be created when the monitoring starts, whereas in online monitoring events for a future time step can still be undefined at a previous time step, but have to be defined later on.


There exists different types of runtime verification algorithms.
The two main categories are rewriting- and automaton-based.